
\chapter*{Expos\'e \texorpdfstring{VII\textsubscript{A}}{VII-A}\\
Infinitesimal Study of Group Schemes}
\phantomsection
\label{sga3:VIIA:title}
\addcontentsline{toc}{chapter}{Expos\'e VII-A. Infinitesimal Study of Group Schemes}
\markboth{EXPOS\'E VII\textsubscript{A}. INFINITESIMAL STUDY OF GROUP SCHEMES}%
  {EXPOS\'E VII\textsubscript{A}. INFINITESIMAL STUDY OF GROUP SCHEMES}

\begin{center}
  \textit{by P. Gabriel}
\end{center}

In Expos\'e II we restricted ourselves to the study of first-order
differential invariants, and we did not address certain phenomena special
to characteristic \(p>0\) or to characteristic zero. Our purpose in
Part~A of this expos\'e is to fill that gap.

Moreover, the infinitesimal study of a group scheme to arbitrary order is
related to that of the associated formal group; the purpose of the second
part of this expos\'e is to present the first definitions and properties
concerning formal groups.

\renewcommand{\thefootnote}{\fnsymbol{footnote}}
\setcounter{footnote}{0}
\section*{A) Differential Operators and Restricted
  \texorpdfstring{\(p\)}{p}-Lie Algebras\footnote{Part~A of the present
  expos\'e was not treated seriously in the oral seminars.}}
\phantomsection
\label{sga3:VIIA:A}
\addcontentsline{toc}{section}{A) Differential Operators and Restricted p-Lie Algebras}
\renewcommand{\thefootnote}{\arabic{footnote}}
\setcounter{footnote}{0}

\section*{1. Differential Operators}
\phantomsection
\label{sga3:VIIA:1}
\addcontentsline{toc}{section}{1. Differential Operators}

In this section, as well as in \SGARefLink{sga3:VIIA:2}{Sections~2} and~3, \(S\) denotes a fixed
scheme, and all products under consideration are cartesian products in
the category of \(S\)-schemes.\footnote{In particular, if \(X\) and \(Y\)
are two \(S\)-schemes, \(X\times_S Y\) is denoted simply by \(X\times Y\).
For the material in \SGARefLink{sga3:VIIA:1}{Sections~1} and~2, one may also consult \emph{Groupes
alg\'ebriques} by Demazure and Gabriel, II.4, nos.~5--6; see also
Jantzen, \emph{Representations of Algebraic Groups}, I.7.}
If \(X\) is an \(S\)-scheme, we write \(p_{X/S}\), \(p_X\), or simply
\(p\) for the structural morphism from \(X\) to \(S\).

\medskip
\noindent\textbf{1.1.}
\phantomsection
\label{sga3:VIIA:1.1}
Let \(u:Y\to X\) be a morphism of \(S\)-schemes, and endow the direct
image \(u_*(\mathcal O_Y)\) of the structure sheaf of \(Y\) with the
\(\mathcal O_X\)-module structure induced by \(u\). The sheaf
\[
  \mathcal H=
  \mathcal Hom_{p_X^{-1}(\mathcal O_S)}
    \bigl(\mathcal O_X,u_*(\mathcal O_Y)\bigr)
\]
of homomorphisms of \(p_X^{-1}(\mathcal O_S)\)-modules from
\(\mathcal O_X\) to \(u_*(\mathcal O_Y)\) is therefore naturally an
\(\mathcal O_X\)-bimodule. If \(U\) is an open subset of \(X\), and
\(f\) and \(d\) are sections over \(U\) of \(\mathcal O_X\) and
\(\mathcal H\), respectively, then \(fd\) and \(df\) are the morphisms
\[
  g\longmapsto f\,d(g)
  \qquad\text{and}\qquad
  g\longmapsto d(fg)
\]
from \(\mathcal O_X\) to \(u_*(\mathcal O_Y)\). From now on we write
\((\operatorname{ad}f)(d)\) in place of \(fd-df\).

\medskip
\noindent\textbf{Definition 1.1.1.}
\phantomsection
\label{sga3:VIIA:1.1.1}
\textit{An \(S\)-deviation of order \(\leq n\) is a pair
\(D=(u,d)\), consisting of a morphism of \(S\)-schemes \(u:Y\to X\)
and a morphism of \(p_X^{-1}(\mathcal O_S)\)-modules
\[
  d:\mathcal O_X\longrightarrow u_*(\mathcal O_Y),
\]
such that, for every open subset \(U\) of \(X\) and every sequence of
\(n+1\) sections \(f_0,\ldots,f_n\in\mathcal O_X(U)\), one has}
\begin{equation}
  (\operatorname{ad}f_0)(\operatorname{ad}f_1)\cdots
  (\operatorname{ad}f_n)(d)=0.
  \tag{\(\ast_n\)}
  \label{sga3:VIIA:eq:star-n}
\end{equation}
\footnote{Equivalently, for every \(x\in X\) and
\(f_0,\ldots,f_n,g\in\mathcal O_{X,x}\),
\[
  (\operatorname{ad}f_0)\cdots(\operatorname{ad}f_n)(d_x)(g)=0.
\]
Recall also that the adjunction isomorphism
\[
 \theta:
 \operatorname{Hom}_{p_X^{-1}(\mathcal O_S)}
   \bigl(\mathcal O_X,u_*(\mathcal O_Y)\bigr)
 \xrightarrow{\sim}
 \operatorname{Hom}_{p_Y^{-1}(\mathcal O_S)}
   \bigl(u^{-1}\mathcal O_X,\mathcal O_Y\bigr)
\]
sends \(d\) to \(d'=\varepsilon\circ u^{-1}(d)\), where
\(\varepsilon:u^{-1}u_*(\mathcal O_Y)\to\mathcal O_Y\) is canonical.
Conversely,
\(\theta^{-1}(d')=u_*(d')\circ\eta\), with
\(\eta:\mathcal O_X\to u_*u^{-1}\mathcal O_X\) canonical. Hence \(d\)
satisfies \((\ast_n)\) if and only if \(d'\) satisfies the analogous
identity for all local sections of \(u^{-1}\mathcal O_X\).}

In this case we also say that \(d\) is an \(S\)-deviation of \(u\) of
order \(\leq n\). In particular, an \(S\)-deviation of \(u\) of order
\(\leq0\) is a morphism of \(\mathcal O_X\)-modules from
\(\mathcal O_X\) to \(u_*(\mathcal O_Y)\), that is, an element of
\(\Gamma(Y,\mathcal O_Y)\).

\medskip
\noindent\textbf{Definition 1.1.2.}
\phantomsection
\label{sga3:VIIA:1.1.2}
\textit{A morphism of \(p_X^{-1}(\mathcal O_S)\)-modules
\[
  d:\mathcal O_X\longrightarrow u_*(\mathcal O_Y)
\]
is an \(S\)-deviation of \(u\) if, for every point \(y\in Y\), there
exist an open neighborhood \(U\) of \(u(y)\) in \(X\) and an open
neighborhood \(V\) of \(y\) in \(Y\) satisfying:
\begin{enumerate}
\item[\textup{a)}] \(u(V)\subset U\);
\item[\textup{b)}] if \(v:V\to U\) is induced by \(u\), there is an
integer \(n\) such that the morphism
\(\mathcal O_U\to v_*(\mathcal O_V)\) induced by \(d\) is an
\(S\)-deviation of \(v\) of order \(\leq n\).
\end{enumerate}}\footnote{If \(X\) and \(u\) are quasi-compact, every
\(S\)-deviation of \(u\) is therefore of order \(\leq n\) for some
integer \(n\).}

If \(d\) is an \(S\)-deviation of \(u\), we also call the pair
\(D=(u,d)\) an \(S\)-deviation, and we may write
\[
  Y\xrightarrow{D}X
  \qquad\text{or}\qquad
  Y\mathop{\longrightarrow}^{d}_{u}X.
\]
When \(d\) is the homomorphism of algebras
\[
  u^\natural:\mathcal O_X\longrightarrow u_*(\mathcal O_Y)
\]
corresponding to the morphism \(u:Y\to X\), we also write \(u\) in
place of \(D\).

\medskip
\noindent\textbf{Remarks 1.1.3.}\footnote{These remarks were added by
the editors because they are useful in \SGARefLink{sga3:VIIA:1.3}{1.3}, \SGARefLink{sga3:VIIA:1.4}{1.4}, and \SGARefLink{sga3:VIIA:2.1}{2.1}.}
\phantomsection
\label{sga3:VIIA:1.1.3}
Let \(\operatorname{Dev}(u)\) (respectively,
\(\operatorname{Dev}_{\leq n}(u)\)) be the set of \(S\)-deviations of
\(u\) (respectively, of order \(\leq n\)). It has a natural
\(\mathcal O_Y(Y)\)-module structure: if
\(\lambda\in\mathcal O_Y(Y)\), then \(\lambda d\) is the deviation
that sends \(f\) to \(\lambda d(f)\) for every local section \(f\) of
\(\mathcal O_X\).

For every open subset \(V\) of \(Y\), set
\(\mathcal{Dev}(u)(V)=\operatorname{Dev}(u|_V)\). Thus
\(\mathcal{Dev}(u)(V)\) is the set of
\begin{align*}
 d_V&\in
 \operatorname{Hom}_{p_X^{-1}(\mathcal O_S)}
   \bigl(\mathcal O_X,(u|_V)_*(\mathcal O_V)\bigr)\\
 &\simeq
 \operatorname{Hom}_{p_Y^{-1}(\mathcal O_S)}
   \bigl((u|_V)^{-1}\mathcal O_X,\mathcal O_V\bigr)\\
 &\simeq
 \operatorname{Hom}_{p_Y^{-1}(\mathcal O_S)}
   \bigl(u^{-1}\mathcal O_X,\mathcal O_Y\bigr)(V)
\end{align*}
such that, for every open \(U\subset X\), the map
\[
 d_V(U):\mathcal O_X(U)\longrightarrow
 \mathcal O_Y\bigl(u^{-1}(U)\cap V\bigr)
\]
satisfies \((\ast_n)\). This defines a presheaf of
\(\mathcal O_Y\)-modules on \(Y\), and it is easily seen to be a sheaf
(more precisely, a subsheaf of
\(\mathcal Hom_{p_Y^{-1}(\mathcal O_S)}
  (u^{-1}\mathcal O_X,\mathcal O_Y)\)).

\medskip
\noindent\textbf{1.2.}
\phantomsection
\label{sga3:VIIA:1.2}
Consider two \(S\)-deviations \(D=(u,d)\) and \(E=(v,e)\):
\[
  Z\mathop{\longrightarrow}^{e}_{v}Y
  \mathop{\longrightarrow}^{d}_{u}X.
\]
As \(U\) ranges over the open subsets of \(X\), the composites
\[
 \Gamma(U,\mathcal O_X)
 \xrightarrow{d(U)}
 \Gamma(u^{-1}U,\mathcal O_Y)
 \xrightarrow{e(u^{-1}U)}
 \Gamma(v^{-1}u^{-1}U,\mathcal O_Z)
\]
define an \(S\)-deviation of \(uv\), denoted by \(de\). If \(d\) has
order \(\leq m\) and \(e\) has order \(\leq n\), then \(de\) has order
\(\leq m+n\). We also write
\begin{equation}
  D\circ E=(uv,de)
  \tag{\(\dagger\)}
  \label{sga3:VIIA:eq:composition}
\end{equation}
\footnote{With this notation, \(de\) means the composite ``\(d\)
followed by \(e\).''}
and call \(D\circ E\), or \(DE\), the \emph{composite
\(S\)-deviation}. When \(d=u^\natural\), so that \(D=u\) under the
convention of \SGARefLink{sga3:VIIA:1.1}{1.1}, we also say that \(DE\) is the \emph{image of \(E\)
under \(u\)}.

The map \((D,E)\mapsto D\circ E\) allows us to speak from now on of the
\emph{category of \(S\)-deviations}, whose objects are \(S\)-schemes
and whose morphisms are \(S\)-deviations.\footnote{Often one considers
only the \(S\)-deviations of \(\operatorname{id}_X\); they form the
algebra of \(S\)-differential operators on \(X\), cf.~\SGARefLink{sga3:VIIA:1.4}{1.4} below. The
more general framework provides a convenient functorial language for
statements such as the following: if \(G\) is an \(S\)-group, the
algebra of left-translation-invariant \(S\)-differential operators on
\(G\) is isomorphic to the algebra of \(S\)-deviations of the unit
section \(\varepsilon:S\to G\); cf.~\SGARefLink{sga3:VIIA:2.1}{2.1} and~\SGARefLink{sga3:VIIA:2.4}{2.4}.}

\medskip
\noindent\textbf{Definition 1.2.0.}\footnote{The editors expanded this
paragraph and assigned the number \SGARefLink{sga3:VIIA:1.2.0}{1.2.0} to this definition and 1.2.1
to the lemma that follows.}
\phantomsection
\label{sga3:VIIA:1.2.0}
\textit{Let \(w:Z\to X\) be an \(S\)-morphism. An
\(S\)-derivation of \(w\), or an \(S\)-derivation of \(\mathcal O_X\)
into \(w_*(\mathcal O_Z)\), is a morphism of
\(p^{-1}(\mathcal O_S)\)-modules
\[
  d:\mathcal O_X\longrightarrow w_*(\mathcal O_Z)
\]
such that, for every open \(U\subset X\) and
\(f,g\in\mathcal O_X(U)\),
\[
  d(fg)=w^\natural(f)d(g)+w^\natural(g)d(f).
\]
Then \(d\) is a deviation of \(w\) of order \(\leq1\), vanishing on
the unit section of \(\mathcal O_X\). We denote the set of
\(S\)-derivations of \(w\) by \(\operatorname{Der}_S(w)\); it is an
\(\mathcal O(Z)\)-module.}

With the notation of \SGARefLink{sga3:VIIA:1.2}{1.2}, take
\[
  Y=I_Z=\operatorname{Spec}\mathcal O_Z[t],\qquad t^2=0,
\]
let \(v\) be the zero section \(\tau:Z\to I_Z\), defined by the
homomorphism \(\mathcal O_Z[t]\to\mathcal O_Z\) sending \(t\) to zero,
and let \(e\) be the morphism of \(\mathcal O_Z\)-modules
\(\sigma:\mathcal O_Z[t]\to\mathcal O_Z\) defined by
\(\sigma(1)=0\) and \(\sigma(t)=1\), which it is convenient to denote
by \(\partial_t\).\footnote{The editors added this notation
\(\partial_t\).}

If \(u:I_Z\to X\) is a morphism satisfying \(w=u\circ\tau\), then
\(\sigma\circ u^\natural\) is an \(S\)-derivation from
\(\mathcal O_X\) to \(w_*(\mathcal O_Z)\). Conversely, an
\(S\)-derivation \(d\) determines the morphism \(u:I_Z\to X\) which
equals \(w\) on the underlying spaces and satisfies
\[
  u^\natural(f)=w^\natural(f)+d(f)t
\]
for every local section \(f\) of \(\mathcal O_X\).

\medskip
\noindent\textbf{Lemma 1.2.1.}
\phantomsection
\label{sga3:VIIA:1.2.1}
\textit{Let \(E=(\tau,\partial_t)\) be the deviation of
\(\tau:Z\to I_Z\) defined above. For every \(S\)-morphism \(w:Z\to X\),
the map \(u\mapsto u\circ E\) is a bijection between the
\(S\)-morphisms \(u:I_Z\to X\) such that \(u\circ\tau=w\), and the
\(S\)-derivations of \(w\).}

\medskip
\noindent\textbf{1.2.2.}
\phantomsection
\label{sga3:VIIA:1.2.2}
Let \(d\) be an \(S\)-deviation of \(u:Y\to X\). On the one hand, \(d\)
is plainly an \(S'\)-deviation of \(u\) for every morphism \(s:S\to S'\).

On the other hand, let \(t:T\to S\) be a morphism with target \(S\),
let \(u_T:Y_T\to X_T\) be obtained from \(u\) by base change, and let
\(t_Y:Y_T\to Y\) and \(t_X:X_T\to X\) be the canonical projections.
There exists a unique \(T\)-deviation of \(u_T\), denoted \(d_T\) or
\(d\times T\), which satisfies
\[
  t_Xd_T=dt_Y
\]
in the sense of \eqref{sga3:VIIA:eq:composition}. Thus, for every open
\(U\subset X\), the following diagram is commutative:
\begin{center}
  \phantomsection
  \label{sga3:VIIA:1.2.2-base-change-square}
  \begin{tikzcd}[column sep=large,row sep=large]
\mathcal O(U)
  \arrow[r,"t_X^{\natural}"]
  \arrow[d,"d(U)"']
&
\mathcal O(U\times T)
  \arrow[d,"d_T(U\times T)"]
\\
\mathcal O(u^{-1}U)
  \arrow[r,"t_Y^{\natural}"]
&
\mathcal O(u^{-1}U\times T)\,.
\end{tikzcd}

\end{center}
\footnote{Explicitly, if \(V\) is an affine open subset of \(S\), and
\(U\) (respectively, \(U'\)) is an affine open subset of \(X\)
(respectively, \(T\)) over \(V\), then
\[
 \mathcal O_{X\times T}(U\times U')
 =\mathcal O_X(U)\otimes_{\mathcal O_S(V)}\mathcal O_T(U'),
\]
and \(d_T(U\times U')\) is the composite
\[
 \mathcal O_X(U)\otimes_{\mathcal O_S(V)}\mathcal O_T(U')
 \xrightarrow{d(U)\otimes\mathrm{id}}
 \mathcal O_Y(u^{-1}U)\otimes_{\mathcal O_S(V)}\mathcal O_T(U')
 \longrightarrow
 \mathcal O_{Y\times T}(u^{-1}U\times U').
\]
The author left to the reader the verification that \(d_T\) is
well-defined, and the editors do the same.}
If \(D=(u,d)\), we also write \(D_T=(u_T,d_T)\), and say that \(d_T\)
and \(D_T\) are obtained from \(d\) and \(D\) by base change.

\medskip
\noindent\textbf{1.2.3.}
\phantomsection
\label{sga3:VIIA:1.2.3}
For example, let \(u:Y\to X\) and \(v:Z\to T\) be two \(S\)-morphisms,
and let \(d\) and \(e\) be \(S\)-deviations of \(u\) and \(v\).
There is a commutative diagram
\begin{center}
  \phantomsection
  \label{sga3:VIIA:1.2.3-product-scheme-square}
  \begin{tikzcd}[column sep=large,row sep=large]
X\times T
&
Y\times T
  \arrow[l,"u_T"']
\\
X\times Z
  \arrow[u,"v_X"]
&
Y\times Z
  \arrow[l,"u_Z"']
  \arrow[u,"v_Y"']
  \arrow[ul,"u\times v"']
\end{tikzcd}

\end{center}
and we denote by \(d\times e\), the \emph{product of \(d\) and \(e\)},
the \(S\)-deviation of \(u\times v\) equal to
\[
  d_Te_Y=e_Xd_Z
\]
under the convention of \eqref{sga3:VIIA:eq:composition}. More
explicitly, for every open \(U\subset X\times T\), set
\[
  W=v_Y^{-1}u_T^{-1}U=u_Z^{-1}v_X^{-1}U.
\]
Then the following diagram is commutative:
\begin{center}
  \phantomsection
  \label{sga3:VIIA:1.2.3-product-deviation-square}
  \begin{tikzcd}[column sep=huge,row sep=large]
\mathcal O(U)
  \arrow[r,"d_T(U)"]
  \arrow[d,"e_X(U)"']
  \arrow[dr,dashed,"(d\times e)(U)" description]
&
\mathcal O(u_T^{-1}U)
  \arrow[d,"e_Y(u_T^{-1}U)"]
\\
\mathcal O(v_X^{-1}U)
  \arrow[r,"d_Z(v_X^{-1}U)"]
&
\mathcal O(W)\,.
\end{tikzcd}

\end{center}
If \(D=(u,d)\) and \(E=(v,e)\), we also write
\[
  D\times E=(u\times v,d\times e).
\]
\footnote{Source note: the French sentence at this point prints
\(E=(v,d)\); the established notation is \(E=(v,e)\).}

